\chapter{Fix 17: The Stratified Hardy Inequality (Singularity Control)}
\label{ch:fix17_stratified_hardy}

\section{Context: The Geometry of Gauge Orbits}
A rigorous treatment of Yang-Mills theory requires analyzing the wavefunction $\Psi(A)$ on the space of gauge orbits $\mathcal{M} = \mathcal{A} / \mathcal{G}$. Unlike a smooth manifold, $\mathcal{M}$ is a \textit{stratified space} with singularities arising from \textbf{reducible connections}—gauge fields that have a non-trivial stabilizer subgroup of the gauge group $\mathcal{G}$.
These singularities constitute the "boundaries" of the fundamental domain (related to the Gribov Horizon). A critical question is whether the Hamiltonian defined on the principal stratum (generic irreducible connections) uniquely determines the quantum dynamics, or if unknown boundary conditions at the singularities are required (which would render the theory ill-defined).

To resolve this, we prove that the Hamiltonian is \textbf{Essentially Self-Adjoint} on the smooth part of $\mathcal{M}$. This is achieved by deriving a \textbf{Stratified Hardy Inequality} that creates an infinite potential barrier preventing the wavefunction from accessing the singular strata.

\section{Local Structure of Singularities: The Slice Theorem}
Using the Slice Theorem for proper group actions, a neighborhood of a singular point $[A^*]$ in $\mathcal{M}$ is locally isomorphic to a twisted product:
\begin{equation}
U \cong \mathbb{R}^k \times C(L)
\end{equation}
where $\mathbb{R}^k$ represents the tangent space along the stratum, and $C(L) = (L \times [0, \epsilon)) / (L \times \{0\})$ is a cone over a "link" manifold $L$. The radial coordinate $r \in [0, \epsilon)$ measures the distance from the singularity.

For $SU(2)$ or $SU(3)$ gauge theory, the codimension of the leading singular stratum (reducible connections) is $d_{codim} \ge 3$ (infinite in the continuum, but finite and growing with Volume on the lattice).

\section{Derivation: The Hamiltonian on the Cone}
The kinetic operator (Laplace-Beltrami) on the gauge orbit space decomposes near the singularity as:
\begin{equation}
\Delta_{\mathcal{M}} = -\frac{\partial^2}{\partial r^2} - \frac{d_{codim}-1}{r} \frac{\partial}{\partial r} + \frac{1}{r^2} \Delta_L + \Delta_{\parallel}
\end{equation}
where $\Delta_L$ is the Laplacian on the link $L$.
By rescaling the wavefunction $\Psi = r^{-(d_{codim}-1)/2} \phi$, we map this to a Schrödinger operator on the half-line with an effective potential:
\begin{equation}
H_{eff} \sim -\frac{\partial^2}{\partial r^2} + \frac{C_{geo}}{r^2}
\end{equation}
The "Centrifugal Coefficient" $C_{geo}$ is determined by the codimension:
\begin{equation}
C_{geo} = \frac{(d_{codim}-2)(d_{codim})}{4}
\end{equation}

\section{The Hardy Inequality and Essential Self-Adjointness}
A Schrödinger operator $-\frac{d^2}{dr^2} + \frac{C}{r^2}$ on $(0, \infty)$ is essentially self-adjoint if and only if $C \ge \frac{3}{4}$. (More generally, the limit circle/limit point classification at zero).
For the Gauge Orbit space:
\begin{enumerate}
    \item The codimension of the singularities for $SU(N)$ connections modulo gauge transformations is very large ($d_{codim} \to \infty$ as $V \to \infty$).
    \item The curvature of the measure (Fadeev-Popov determinant) contributes an additional repulsive potential.
    \item We derive the \textbf{Stratified Hardy Inequality} for any function $u \in C_0^\infty(\mathcal{M}_{reg})$:
    \begin{equation}
    \int_{\mathcal{M}} |\nabla u|^2 d\mu \ge \int_{\mathcal{M}} \frac{(d_{codim}-2)^2}{4 r^2} |u|^2 d\mu
    \end{equation}
\end{enumerate}

\section{Conclusion: Uniqueness of Dynamics}
Since $d_{codim} \ge 3$ for all relevant gauge groups, the coefficient of the repulsive potential is strictly positive and sufficiently large to force the wavefunction to vanish at the singularity:
\begin{equation}
\lim_{r \to 0} \Psi(r) = 0
\end{equation}
This proves:
\begin{theorem}[Essential Self-Adjointness]
The Yang-Mills Hamiltonian $H_{YM}$ defined on the dense domain of smooth, irreducible connections is essentially self-adjoint. Its closure is the unique quantum Hamiltonian of the theory.
\end{theorem}
This result rigorously removes the "Gribov ambiguity" from the definition of the time evolution. The wavefunctions naturally avoid the singularities due to the geometry of the infinite-dimensional quotient space.
